How can you measure curvature without stepping off a surface? Why does a
loop through a torus behave differently from a loop on a sphere? How does
a uniform choice of unitary matrix produce a precise law of eigenvalue
repulsion?

This book develops the geometry behind these questions. Coordinates lead
to smooth manifolds, vector fields and differential forms. Stokes' theorem
connects differentiation with integration; Gauss--Bonnet connects curvature
with topology. Fundamental groups, covering spaces and homology distinguish
spaces that look the same locally.

Matrix Lie groups bring rotations and quaternions into the picture.
Haar measure and Weyl's formula describe random unitary matrices.
Hyperbolic geometry connects matrix motion with continued fractions.
The Riemann mapping theorem gives holomorphic maps onto the disc, and
Crofton's formula measures curves by counting crossings with lines.
Eleven appendices develop the supporting geometry, from inversion and
conics to cubic curves, sphere packings and convex bodies.

Proofs explain why each step is needed. Advanced results quoted without
proof are identified. Every exercise has a hint and a complete solution,
with the idea of the argument set out before the writeup.
